\section{Analysis Layer~2 --- Design Values}
\label{sec:ch5-layer2}

This section presents seven design values developed interpretively through reflexive thematic analysis~\cite{braun2023thematic} of expert discussions. It captures value orientations inferred from how experts reasoned about what biophilic interaction should enable, constrain, or leave indeterminate when translating outdoor experience indoors (\textbf{RQ2}).


\subsection{Sensory Ecology}
Sensory Ecology captures biophilic experience as emerging from multiple, loosely coupled sensory channels, where no single cue carries the full burden of ``being biophilic''. It supports variation across occupants and moments, resisting designs that rely on one modality to stand in for nature. Importantly, it does not imply intensity or immersion. Experts cautioned that tightly synchronised multisensory cues (e.g. aligning sound, light, and movement) risk collapsing variation into overly directive experiences. This orientation appeared in a corridor concept where experts combined airflow variation, sound, and light reflections without designating any cue as primary. They noted that some occupants might notice the shifting sound field, while others sensitive to bodily cues might register the air movement as thermal sensation, and some may be attracted to the reflections. One expert explicitly linked this to neurodivergent perception, and that distributing biophilic qualities across modalities allowed occupants to attune to what felt accessible rather than being directed toward a prescribed sensory experience. 

\begin{figure*}[!ht]
    \centering
    \includegraphics[width=\linewidth]{images/dis26-functionalaxis-alt.pdf}
    \caption{Overview of design values developed from patterns in expert reasoning about what biophilic interaction should enable, constrain, or leave indeterminate. The operational–relational spectrum is an analytic device introduced by the authors and does not imply hierarchy. Values at the operational end frame biophilic qualities as background conditions (sensory ecology, attentional ease, proprioceptive engagement), while values at the relational end frame environments as co-present and partially autonomous (unknowability, non-reciprocity, collective liability).}
    \Description{An overview diagram presenting seven recurring design values identified in expert reasoning about biophilic interaction design. The non-hierarchical spectrum shows values ranging from background conditions supporting bodily inhabitation to relational qualities emphasising co-presence and partial autonomy of the environment.}
    \label{fig:layer2-values}
\end{figure*}

\subsection{Attentional Ease}
Attentional Ease prioritises conditions that allow (or encourage) attention to wander without instruction, while resisting designs that frame attention as something to be captured, directed, or managed. This value surfaced in corridors between meeting rooms, where experts discussed semi-translucent surfaces and low-resolution LED light patterns that could be sensed but not clearly read. These conditions allow attention to momentarily lose focus while walking, creating a brief interval where occupants could pass through without orienting toward information, destination, or task. Because nothing in the space requires recognition or response, it supports mind wandering as a by-product of movement rather than as a designated pause or intervention.

\subsection{Proprioceptive Engagement}
Proprioceptive Engagement captures the importance of bodily movement as a primary source of spatial awareness and environmental understanding. While Attentional Ease prioritises mind drift, Proprioceptive Engagement describes how it invites the body to actively work, to balance on uneven surfaces, adjust posture continuously, negotiate varied terrain, and build spatial knowledge through kinaesthetic feedback. This value prioritises movement as experiential rather than instrumental. Experts (inspired by Ai Weiwei’s Sunflower Seeds installation\footnote{\url{https://www.tate.org.uk/whats-on/tate-modern/unilever-series/unilever-series-ai-weiwei-sunflower-seeds}}) sketched an atrium floor surface composed of small, loose elements that shifted underfoot. Rather than functioning as a path or obstacle, the surface was intended to register through balance, sound, and gait adjustments. Experts emphasised that occupants could choose to walk across, skirt its edge, or ignore it entirely, and allowing biophilic experience to arise through doing rather than observation.

\subsection{Irreversibility}
Irreversibility is the recognition that certain biophilic experiences derive meaning from processes that cannot be undone, reset, or perfectly controlled. This value prioritises environments that register consequence, protecting traces of use, accumulation, and change, while resisting design logics that prioritise constant reset and reversibility. Experts proposed an office soundscape that gradually incorporated sounds generated by previous occupants’ movements, described as ``carrying the memory of who came here before.'' In another, a soft space divider was designed to deform in response to airflow each time a door opened, so that over the course of a day it registered patterns of arrival and presence as a slowly changing form. In both cases, traces accumulated through use, serving as a record of inhabitation.

\subsection{Unknowability}
Unknowability is the value placed on environments that cannot be fully predicted, explained, or mastered. This value prioritises ambiguity, variation, and partial legibility. It protects curiosity and ongoing interpretation of occupants, while resisting design logics that privilege predictability, optimisation, or complete transparency. Experts repeatedly described nature as resisting complete understanding, through changing conditions, partial visibility, or sensory ambiguity, and framed this unpredictability as engaging rather than problematic. Experts described ambient light and sound that shifted in response to outdoor conditions and occupancy patterns, without exposing a control logic. They noted that occupants could sense recurring tendencies, such as the room feeling different on overcast days or later in the afternoon, without being able to fully anticipate or explain these changes.

\subsection{Non-Reciprocity}
Non-Reciprocity captures the idea that interactive biophilic experiences do not necessarily have to respond to human presence. This value prioritises persistence over feedback. It protects a sense of scale, autonomy, and indifference associated with natural environments, while resisting design logics that equate interaction with responsiveness (e.g. smart building systems that are designed to ``notice'' or ``adapt''~\cite{profita2015wall}). Experts explored an office waiting area where a biophilic installation combined living materials with (computationally driven) mechanical processes that evolved continuously over time. Humidity levels and material movement followed internal system logic rather than responding to occupancy or individual presence. People entering the space encountered the installation mid-process, sometimes quiet, sometimes visibly active, but nothing acknowledged their arrival or departure.

\subsection{Collective Liability}
Collective Liability captures biophilic experience where occupants are implicated together, rather than individually insulated from. This value foregrounds shared exposure and consequence, resisting design logics that buffer people from one another through personalisation, optimisation, or individual control. Biophilic environments are lived with as common conditions that require mutual accommodation. Experts reflected on elevator lobbies where collective patterns of use produced perceptible environmental consequences (also explored by ~\citet{loh2010please}). They described how, over the course of a morning, frequent elevator use led the space to gradually retain warmth, reflecting the cumulative energy consumed. As more people chose the elevator over the stairs, the lobby increasingly carried this residual heat, encountered by everyone passing through. No individual action caused the change; instead, the atmosphere registered the shared impact of collective choice as a common condition rather than a personalised response.

\subsection{Summary}
This layer surfaces biophilic interaction design values constructed through analysis of experts’ reasoning across the design sessions. We visualise these values along a spectrum from operational to relational in \autoref{fig:layer2-values}. This spectrum is author-developed interpretive device to organise the values. It reflects a recurring distinction between values that frame biophilic interaction as supporting inhabitation through background conditions (e.g. sensory ecology, attentional ease, proprioceptive engagement), and values that frame it as ongoing engagement with partially autonomous, evolving environmental conditions (e.g. unknowability, non-reciprocity, collective liability). The spectrum clarifies how different orientations shift the role of interaction, rather than defining a fixed categorisation. At the operational end, values orient biophilic features toward supporting everyday inhabitation---for example through movement, attentional ease, or bodily accommodation---without requiring interpretation or overt engagement. Toward the relational end, values foreground how biophilic environments are encountered and lived with over time, emphasising asymmetry, ambiguity, persistence, and shared exposure rather than responsiveness or individual optimisation. Collectively, they mark recurring departures from dominant assumptions in smart building and interaction design, such as optimisation, predictability, and individual control (\autoref{subsec:ch5-layer2-discussion}). 

